\section{Forced kinetic amplification at nearly critical mass}
\label{sec:r12-amplification}

We return to the pure mass-critical equation
\[
 iu_t=\Delta^2u-|u|^{8/N}u,\qquad
 p=2+8/N,\qquad
 M=\|u\|_2^2,\quad A=\|\Delta u\|_2^2,\quad
 E=\frac A2-\frac{\|u\|_p^p}{p}.
\]
Let $M_c$ be the mass in the sharp interpolation inequality
\[
 \|f\|_p^p\le\frac p2\left(\frac{M(f)}{M_c}\right)^{4/N}A(f).
\]
Solutions are strong conservative $H^2$ solutions for which the integrated
localized mass and virial identities hold by approximation. No spatial
moment or radial symmetry is imposed.

\begin{theorem}[Quantitative kinetic amplification]
\label{thm:r12-amplification}
There are $d_0,C>0$, depending only on the scalar sharp interpolation
problem, with the following property. If a forward-global solution has
\[
 P=0,\qquad E=-e<0,\qquad
 d=\left(\frac M{M_c}\right)^{4/N}-1\in(0,d_0),\qquad
 B=\sup_{t\ge0}A(u(t))<\infty,
\]
then
\begin{equation}\label{eq:r12-amplification}
 e\le C d^{5/3}B.
\end{equation}
Consequently, in any maximal strong $H^2$ flow with the usual continuation
alternative,
\begin{equation}\label{eq:r12-maximal-amplification}
 \sup_{0\le t<T_+}A(u(t))\ge C^{-1}|E|d^{-5/3}.
\end{equation}
The latter statement allows finite-time blowup, global unbounded motion,
and bounded global motion.
\end{theorem}

The instantaneous sharp interpolation inequality gives only
$A(t)\ge2|E|/d$. Thus \eqref{eq:r12-maximal-amplification} forces a
larger kinetic scale, by a factor comparable to $d^{-2/3}$.
The result does not supply an upper bound on that scale and therefore
does not settle finite-time blowup or growup for all such initial data.

\subsection{A quantitative constant-phase estimate}
For a nonzero scalar complex field define
\[
 \mathfrak d(f)=1-
 \frac{\|f\|_p^p}{(p/2)(M(f)/M_c)^{4/N}A(f)}.
\]

\begin{lemma}[Distance to a real field]
\label{lem:r12-realphase}
There is $C<\infty$ such that every nonzero $f\in H^2(\mathbb R^N;\mathbb C)$
admits a constant phase $\theta$ for which
\begin{equation}\label{eq:r12-realphase}
 \frac{M(\operatorname{Im}(e^{-i\theta}f))}{M(f)}+
 \frac{A(\operatorname{Im}(e^{-i\theta}f))}{A(f)}
 \le C\mathfrak d(f).
\end{equation}
No uniqueness of real optimizers or nondegeneracy of the real linearized
operator is assumed.
\end{lemma}
\begin{proof}
Amplitude and critical spatial scaling reduce the proof to $M(f)=M_c$
and $A(f)=1$. If the result failed, since the left side is at most two,
there would be a sequence with $\mathfrak d(f_n)\to0$ for which the ratio
of the right defect to the best left side tends to zero. Sharp-optimizer
compactness supplies, after translations and phases,
$f_n\to Q$ strongly in $H^2$, for a real normalized optimizer $Q$.
Choose the remaining phase so that $f_n=X_n+iY_n$ and
$\langle Q,Y_n\rangle=0$.
Put $r=4/N$, $\kappa=p/2=1+r$, $a_n=A(Y_n)$ and $y_n=M(Y_n)$.
The transverse Taylor expansion, with the Sobolev remainder estimate,
gives
\[
 \|X_n+iY_n\|_p^p
 =\|X_n\|_p^p+\kappa\int |X_n|^{p-2}Y_n^2
       +o(\|Y_n\|_{H^2}^2).
\]
For $2<p\le4$ the remainder is bounded by $C\|Y_n\|_p^p$;
for $p>4$ use the sum of the fourth-order mixed term and
$C\|Y_n\|_p^p$. Sharp interpolation for $X_n$ yields
\[
 \|X_n\|_p^p\le\kappa(1-y_n/M_c)^r(1-a_n).
\]
Potential form convergence therefore implies
\[
 \kappa-\|f_n\|_p^p
 \ge\kappa\langle Y_n,L_-Y_n\rangle-o(\|Y_n\|_{H^2}^2),
 \qquad L_-=\Delta^2+r/M_c-|Q|^{p-2}.
\]
The sharp phase form has kernel exactly $\mathbb RQ$ and is coercive
on $Q^\perp$, as follows from density strict convexity and the sharp
interpolation inequality. Hence the defect controls
$\|Y_n\|_{H^2}^2$, a contradiction. Compactness of normalized optimizers
makes the coercivity constant uniform. This proves the lemma.
\end{proof}

This is a real two-component polarization statement. Treating $f$ as a
single complex internal component would give a vacuous rank-one assertion
and would not prove \eqref{eq:r12-realphase}.

\subsection{A slowly moving concentration center}
Critical scaling allows us to assume $B=1$. Throughout the proof $M$ then
lies in a fixed compact interval about $M_c$, and
\begin{equation}\label{eq:r12-Arange}
 0<2e/d\le A(t)\le1.
\end{equation}
Since
\[
 \mathfrak d(u(t))=\frac{d-2e/A(t)}{1+d}\le d,
\]
Lemma~\ref{lem:r12-realphase} gives, separately at each time, a phase
such that $u=e^{i\theta}(X+iY)$ with
\begin{equation}\label{eq:r12-phase-small}
 \|Y\|_{H^2}\le C\sqrt d,\qquad \|X\|_{H^2}\le C.
\end{equation}
No time derivative of the phase is used.

For a real smooth bounded test function $\phi$, integration by parts in
the mass identity gives
\begin{equation}\label{eq:r12-density-id}
 \frac d{dt}\int\phi|u|^2
 =2\operatorname{Im}\int
 \bigl(2\nabla\phi\cdot\nabla\overline u+(\Delta\phi)\overline u\bigr)
             \Delta u.
\end{equation}
The local nonlinearity cancels. Substituting
\eqref{eq:r12-phase-small} proves
\[
 \left|\frac d{dt}\int\phi|u|^2\right|
 \le C\sqrt d\bigl(\|\nabla\phi\|_\infty+\|\Delta\phi\|_\infty\bigr).
\]
In particular, for $R\ge1$ and fixed $b$,
\begin{equation}\label{eq:r12-density-speed}
 \left|\frac d{dt}\int\chi((x-b)/R)|u|^2\right|
 \le C\sqrt d/R.
\end{equation}
These statements hold in integrated form for the stated solution class.

Choose $\varepsilon_0=M_c/32$. For sufficiently small $d_0$, compactness
of sharp near-optimizers and \eqref{eq:r12-Arange} give a radius
$R_0\ge1$ and a choice $c(t)$ satisfying
\begin{equation}\label{eq:r12-core-mass}
 \int_{|x-c(t)|\le R_0}|u(t,x)|^2\,dx\ge M_c-\varepsilon_0.
\end{equation}
The radius may depend on $e/d$; no uniform lower kinetic scale is needed.
Take $\chi=1$ on the unit ball and $\chi=0$ outside the radius-two ball.
If $|c(t)-c(s)|>4R_0$, the integral in
\eqref{eq:r12-density-speed}, centered at $c(s)$, changes by at least
$M_c-3\varepsilon_0$, provided $M-M_c\le\varepsilon_0$.
Thus, for a fixed $\tau>0$,
\[
 |t-s|\le\tau R_0/\sqrt d\quad\Longrightarrow\quad
 |c(t)-c(s)|\le4R_0.
\]
Interpolate $c$ linearly at time spacing $\tau R_0/\sqrt d$. We obtain
one Lipschitz center $a$ satisfying
\begin{equation}\label{eq:r12-center-speed}
 \|a'\|_\infty\le L=C\sqrt d,\qquad
 \sup_t|a(t)-c(t)|\le8R_0.
\end{equation}
Crucially, its speed bound is independent of $R_0$ and $e/d$.
A fixed ball about $a(t)$ retains the core mass in
\eqref{eq:r12-core-mass}.

\begin{lemma}[Uniform mass outside the core]
\label{lem:r12-tailmass}
For every $\eta>0$ there is $R_\eta<\infty$ such that
\begin{equation}\label{eq:r12-tailmass}
 \sup_{t\ge0}\int_{|x-a(t)|>R_\eta}|u(t,x)|^2\,dx
 \le M-M_c+\eta.
\end{equation}
\end{lemma}
\begin{proof}
Otherwise choose $t_k$ and $R_k\to\infty$ violating the inequality.
After translation by $a(t_k)$, a subsequence converges weakly in $H^2$
and strongly locally in $L^2$ to $v$. The fixed core gives
$M(v)\ge M_c-\varepsilon_0$, while the assumed exterior mass implies
$M(v)\le M_c-\eta$. The remainder $w_k$ therefore has limiting mass
at most $M-M_c+\varepsilon_0<M_c$. Quadratic splitting and the
Brezis--Lieb identity give
\[
 -e=E(v)+E(w_k)+o(1).
\]
Both terms on the right are nonnegative by sharp interpolation, since
both masses are strictly subcritical. This is a contradiction.
\end{proof}

The lemma permits a noncompact mass tail of size $M-M_c$. It establishes
neither mass compactness nor kinetic compactness of the full trajectory.

\subsection{Negative exterior kinetic energy in the virial derivative}
Choose a real radial $\phi_R$ equal to $|x|^2/2$ on $|x|\le R$ and
constant on $|x|\ge2R$, with
\[
 D^2\phi_R\le I,\qquad
 \|D^j\phi_R\|_\infty\le C_jR^{2-j}.
\]
One construction prescribes $\phi_R'(r)=rh(r/R)$ with smooth
nonincreasing $h$, equal to one on $[0,1]$ and zero on $[2,\infty)$.
Set
\begin{align*}
 J_R(t)&=2\operatorname{Im}\int
 \nabla\phi_R(x-a(t))\cdot\overline u\nabla u,\\
 K_R(t)&=\int_{|x-a(t)|>2R}|D^2u|^2,\\
 H_R(t)&=\int_{R/2<|x-a(t)|<2R}|D^2u|^2.
\end{align*}
We use the full Hessian in these tails. Its total squared norm equals
$\|\Delta u\|_2^2$, but the two densities are not pointwise identical.

The exact localized quartic identity is
\begin{align*}
 J_R'={}&8\operatorname{Re}\int\partial_{jk}\phi_R\,
           \partial_{ij}\overline u\,\partial_{ik}u
 -4\operatorname{Re}\int\partial_{jk}\Delta\phi_R\,
           \partial_j\overline u\,\partial_ku\\
 &-2\int\Delta^2\phi_R|\nabla u|^2
   +\int\Delta^3\phi_R|u|^2
   -\frac{16}{Np}\int\Delta\phi_R|u|^p
   -2a'\cdot\int D^2\phi_R\operatorname{Im}(\overline u\nabla u).
\end{align*}
The cutoff is centered at $a(t)$. Relative to the global kinetic term
$8A$, the first term has a nonpositive defect everywhere, and its defect
is exactly $-8|D^2u|^2$ outside $2R$. Higher cutoff derivatives contribute
$O(R^{-2})$. Subtracting the total momentum from the last term and using
$P=0$ therefore gives
\begin{equation}\label{eq:r12-retained-virial}
 J_R'\le-16e-8K_R+C\int_{|x-a|>R}|u|^p
          +CR^{-2}+CL\int_{|x-a|>R}|u||\nabla u|.
\end{equation}
Also $|J_R|\le CR$.

Fix small $\eta>0$, put $\mu=M-M_c+\eta$, and take $R\ge2R_\eta$.
For a cutoff $\zeta_R$ vanishing inside $R/2$ and equal to one outside
$R$, Lemma~\ref{lem:r12-tailmass} and the product rule give
\[
 M(\zeta_Ru)\le\mu,\qquad
 A(\zeta_Ru)\le C(K_R+H_R+R^{-2}).
\]
Sharp interpolation and
$\|\nabla f\|_2^2\le\|f\|_2\|\Delta f\|_2$ imply
\begin{align*}
 \int_{|x-a|>R}|u|^p
   &\le C\mu^{4/N}(K_R+H_R+R^{-2}),\\
 \int_{|x-a|>R}|u||\nabla u|
   &\le C\mu^{3/4}\bigl((K_R+H_R)^{1/4}+R^{-1/2}\bigr).
\end{align*}
Choose $d_0$ and $\eta$ so small that $C\mu^{4/N}\le1$.
Absorb the terms containing $K_R$ into its negative coefficient in
\eqref{eq:r12-retained-virial}; for the current term use Young's inequality.
We obtain
\begin{equation}\label{eq:r12-annular-virial}
 J_R'\le-16e+C\mu^{4/N}H_R+CL\mu^{3/4}H_R^{1/4}
             +CL^{4/3}\mu+o_R(1),
\end{equation}
where $o_R(1)\to0$ uniformly in time. No favorable lower bound on the
kinetic defect in the transition annulus was used.

\subsection{Averaging disjoint annuli}
Fix a positive integer $K$ and $R_*\ge2R_\eta$. Let
$R_j=8^jR_*$ for $0\le j<K$. Their transition annuli are disjoint about
the same center, so
\[
 \sum_{j=0}^{K-1}H_{R_j}(t)\le A(t)\le1.
\]
For each $T$ choose an index whose time-averaged $H_{R_j}$ is at most
$K^{-1}$. Integrate \eqref{eq:r12-annular-virial}, use Jensen's inequality
for the fourth root, and bound both endpoint values of $J_{R_j}$:
\[
 16e\le\frac{2C8^{K-1}R_*}{T}
       +\frac{C\mu^{4/N}}K+CL\mu^{3/4}K^{-1/4}
       +CL^{4/3}\mu+\sup_{R\ge R_*}|o_R(1)|.
\]
First let $T\to\infty$ with $K,R_*$ fixed, then $K\to\infty$,
then $R_*\to\infty$, and finally $\eta\downarrow0$. This proves
\[
 e\le CL^{4/3}(M-M_c).
\]
The radius is fixed during each time integration. By
\eqref{eq:r12-center-speed}, $L\le C\sqrt d$, and $M-M_c\le Cd$.
Thus $e\le Cd^{5/3}$ when $B=1$. Undoing critical scaling proves
Theorem~\ref{thm:r12-amplification}.

\subsection{Finite-time excursions from every sharp standing wave}
\begin{corollary}[Standing-wave instability without real-sector nondegeneracy]
\label{cor:r12-standing-instability}
Let $Q$ be any real sharp optimizer normalized by $M(Q)=M_c$,
with its associated standing-wave frequency. Assume a strong $H^2$ local flow and the continuation alternative.
For
\[
 u_d(0)=(1+d)^{N/8}Q,\qquad d>0,
\]
and every prescribed $K>1$, all sufficiently small $d$ admit a finite
forward time $t_d<T_+(d)$ for which
\[
 A(u_d(t_d))>K A(u_d(0)).
\]
In particular the standing wave is orbitally unstable in $H^2$ modulo
constant phase and translations. No radial or $L_+$ nondegeneracy
hypothesis is needed.
\end{corollary}
\begin{proof}
These initial data have $P=0$, $M=(1+d)^{N/4}M_c$ and
$E=-dA(u_d(0))/2$. Theorem~\ref{thm:r12-amplification} gives
\[
 \sup_{0\le t<T_+}\frac{A(u_d(t))}{A(u_d(0))}
 \ge(2C)^{-1}d^{-2/3}.
\]
For small $d$ this is strictly greater than $K$, so a finite time has
the claimed property. Since $u_d(0)\to Q$ in $H^2$, take, for example,
$K=4$ and use
\[
 \inf_{\theta,y}\|u_d(t_d)-e^{i\theta}Q(\cdot-y)\|_{H^2}
 \ge\|\Delta u_d(t_d)\|_2-\|\Delta Q\|_2
 \ge\|\Delta Q\|_2.
\]
This proves orbital instability.
\end{proof}

For $N=1,2,3$ the required strong local flow follows from the $H^2$
algebra and composition estimates, with conservation and the integrated
identities obtained by smooth approximation. In larger dimensions the statement
uses the appropriate strong energy-space local theory. The conclusion
is a finite-time excursion for the perturbed data; its proof does not
identify which of finite blowup, global unbounded evolution, or bounded
large motion occurs.
